\documentclass[11pt]{article}
\usepackage[margin=1in]{geometry}
\usepackage{enumitem}
% =============================================================================
% Preambles/header.tex — shared LaTeX/Beamer preamble for this project
%
% Use from a Beamer lecture by prepending:
%     \input{header}
% and compiling with TEXINPUTS=../Preambles:$TEXINPUTS (the /compile-latex
% skill does this automatically).
%
% PALETTE CONTRACT. Color names defined here MUST match the names in
% Quarto/theme-template.scss so Beamer and Quarto renderings use the same
% palette. The scripts/check-palette-sync.sh script enforces this.
%
% To customize: edit the HEX values below AND the matching $variable in
% Quarto/theme-template.scss, then run scripts/check-palette-sync.sh.
% =============================================================================

% -----------------------------------------------------------------------------
% Engine detection + encoding
%
% Our /compile-latex skill uses XeLaTeX, where inputenc is unnecessary and
% can produce encoding warnings. Only load inputenc under pdfTeX.
% -----------------------------------------------------------------------------
\usepackage{iftex}
\ifPDFTeX
  \usepackage[utf8]{inputenc}
\fi

\usepackage{amsmath, amssymb, amsthm}
\usepackage{booktabs}
\usepackage{graphicx}

% -----------------------------------------------------------------------------
% PALETTE — mirrors Quarto/theme-template.scss variable names.
% Load xcolor and define colors BEFORE anything that references them
% (notably hyperref, hypersetup, and the Beamer color assignments).
% -----------------------------------------------------------------------------
\usepackage{xcolor}

\definecolor{primary-blue}{HTML}{012169}      % headings, accents
\definecolor{primary-gold}{HTML}{B9975B}      % emphasis, borders
\definecolor{highlight-yellow}{HTML}{F2A900}  % markers, alerts
\definecolor{light-bg}{HTML}{E8EDF5}          % title / section backgrounds
\definecolor{jet}{HTML}{1A1A1A}               % body text

% Semantic colors (used in diagrams and annotations)
\definecolor{positive}{HTML}{15803D}          % good / observed / identified
\definecolor{negative}{HTML}{B91C1C}          % bad / problematic / bias
\definecolor{neutral}{HTML}{525252}           % reference / context

% Highlight accents (match .hi-* classes in the SCSS)
\definecolor{hi-slate}{HTML}{314F4F}
\definecolor{hi-green}{HTML}{2E8B57}
\definecolor{hi-red}{HTML}{C41E3A}

% Semantic aliases so skill templates and snippets can write
% \textcolor{accent}{...} without hard-coding "primary-blue".
\colorlet{accent}{primary-blue}
\colorlet{accent2}{primary-gold}

% -----------------------------------------------------------------------------
% Hyperref — loaded AFTER xcolor + palette so link colors resolve.
% -----------------------------------------------------------------------------
\usepackage{hyperref}
\hypersetup{colorlinks=true, linkcolor=primary-blue, urlcolor=primary-blue,
            citecolor=primary-blue, pdfborder={0 0 0}}

% -----------------------------------------------------------------------------
% Course code — set once per deck (after \input{header}, before \begin{document})
% via \coursecode{CS401}. Shown in the footer on every slide so a deck's course
% is identifiable without opening the title page. Empty by default.
% -----------------------------------------------------------------------------
\makeatletter
\newcommand{\coursecode}[1]{\gdef\@coursecodetext{#1}}
\gdef\@coursecodetext{}
\makeatother

% -----------------------------------------------------------------------------
% Beamer theme assignments (only apply under Beamer).
%
% We detect Beamer via \@ifclassloaded{beamer}, which is the documented,
% non-internal way. This must be wrapped in \makeatletter/\makeatother
% because \@ifclassloaded contains an @-catcode character.
% -----------------------------------------------------------------------------
\makeatletter
\@ifclassloaded{beamer}{%
  \usefonttheme{professionalfonts}
  \setbeamercolor{structure}{fg=primary-blue}
  \setbeamercolor{frametitle}{fg=primary-blue}
  \setbeamercolor{title}{fg=primary-blue}
  \setbeamercolor{subtitle}{fg=primary-gold}
  \setbeamercolor{author}{fg=primary-blue}
  \setbeamercolor{institute}{fg=neutral}
  \setbeamercolor{date}{fg=primary-gold}
  \setbeamercolor{itemize item}{fg=primary-blue}
  \setbeamercolor{itemize subitem}{fg=primary-gold}
  \setbeamercolor{alerted text}{fg=primary-gold}
  \setbeamercolor{block title}{fg=white, bg=primary-blue}
  \setbeamercolor{block body}{bg=light-bg}
  % Semantic block variants: block=definition/concept (blue), exampleblock=
  % worked example (green/positive), alertblock=key takeaway or caution (gold).
  % Keep to <=2 colored boxes per slide across ALL three variants (INV-7).
  \setbeamercolor{block title example}{fg=white, bg=positive}
  \setbeamercolor{block body example}{bg=light-bg}
  \setbeamercolor{block title alerted}{fg=white, bg=primary-gold}
  \setbeamercolor{block body alerted}{bg=light-bg}

  % Minimal footer: course code (left) + page X / N (right), no navigation clutter
  \setbeamertemplate{navigation symbols}{}
  \setbeamertemplate{footline}{%
    \color{neutral}\footnotesize\hspace{0.7em}\@coursecodetext\hfill
    \insertframenumber/\inserttotalframenumber\hspace{0.7em}\vspace{0.3em}}
}{}
\makeatother

% -----------------------------------------------------------------------------
% TikZ setup — libraries the snippet gallery uses
% -----------------------------------------------------------------------------
\usepackage{tikz}
\usetikzlibrary{arrows.meta, positioning, calc, decorations.pathreplacing,
                fit, shapes.geometric, backgrounds}

% Shared TikZ styles — snippets and hand-written diagrams can pick these up
% via `\begin{tikzpicture}[every node/.style={...}]` or by naming specific
% styles (e.g., `\node[dag-node] (X) at (0,0) {X};`).
\tikzset{
  % Node styles for causal DAGs and similar
  dag-node/.style = {
    draw=primary-blue, fill=light-bg, rounded corners=2pt,
    minimum width=1.2cm, minimum height=0.9cm, align=center, font=\small
  },
  decision-node/.style = {
    draw=primary-gold, fill=white, diamond, aspect=1.8,
    minimum width=1.8cm, minimum height=1.2cm, align=center, font=\small,
    inner sep=1pt
  },
  % Edge styles — observed vs counterfactual
  observed-edge/.style = {-{Latex[length=2.5mm]}, thick, draw=jet},
  counterfactual-edge/.style = {-{Latex[length=2.5mm]}, thick, draw=neutral, dashed},
  confound-edge/.style = {-{Latex[length=2.5mm]}, thick, draw=negative, dashed},
  % Data-point styles for plots
  observed-dot/.style = {fill=primary-blue, draw=primary-blue, circle, inner sep=1.2pt},
  counterfactual-dot/.style = {fill=white, draw=neutral, circle, inner sep=1.2pt}
}

% -----------------------------------------------------------------------------
% Convenience macros
% -----------------------------------------------------------------------------
\newcommand{\muted}[1]{\textcolor{neutral}{#1}}
\newcommand{\key}[1]{\textcolor{primary-gold}{\textbf{#1}}}
\newcommand{\good}[1]{\textcolor{positive}{#1}}
\newcommand{\bad}[1]{\textcolor{negative}{#1}}

% Standout frame macro: full-bleed dark background for section transitions.
% Beamer-only (uses \begin{frame}/\setbeamercolor/\usebeamerfont) -- do not
% call from an article-class file (e.g. Notes/<CODE>/*-notes.tex). \muted,
% \key, \good, \bad, and \coursecode above are plain xcolor/gdef and are
% safe to use from either class; verified when Notes/article-class support
% was added.
% Expand: \transitionslide{Your transition title}
\newcommand{\transitionslide}[1]{%
  \begingroup
  \setbeamercolor{background canvas}{bg=primary-blue}
  \begin{frame}[plain]
    \centering
    \vfill
    {\usebeamerfont{title}\color{white}\Large #1\par}
    \vfill
  \end{frame}
  \endgroup
}

% Section divider frame (Beamer-only), an alternative to \transitionslide with a
% darker two-line style: near-black (jet) background, a small white label line
% above a larger gold title, and a thin gold rule along the bottom edge.
% Expand: \sectiondivider{Section 2}{Pointers}
\newcommand{\sectiondivider}[2]{%
  \begingroup
  \setbeamercolor{background canvas}{bg=jet}
  \begin{frame}[plain]
    \vspace*{3.0cm}
    {\color{white}\ttfamily\large #1\par}
    \vspace{0.5cm}
    {\color{primary-gold}\LARGE #2\par}
    \begin{tikzpicture}[remember picture, overlay]
      \draw[primary-gold, line width=2.5pt]
        ($(current page.south west)+(0,0.1)$) -- ($(current page.south east)+(0,0.1)$);
    \end{tikzpicture}
  \end{frame}
  \endgroup
}

% -----------------------------------------------------------------------------
% Channel watermark — "Concepts That Click" (YouTube).
%
% Article-class files (CompetitiveExam Q/A sets, Notes, Assignments) get a
% light diagonal draft watermark on every page plus a centered footer naming
% the channel. Beamer decks get a subtle TikZ background watermark instead
% (the footline already carries the course code). Centralized here so every
% MCQ PDF is branded without per-file edits.
% -----------------------------------------------------------------------------
\makeatletter
\@ifclassloaded{article}{%
  \usepackage{draftwatermark}
  \SetWatermarkText{Concepts That Click}
  \SetWatermarkScale{1}
  \SetWatermarkFontSize{2cm}
  \SetWatermarkLightness{0.86}
  \SetWatermarkAngle{45}
  \usepackage{fancyhdr}
  \pagestyle{fancy}
  \fancyhf{}
  \fancyfoot[C]{\footnotesize\textcolor{neutral}{Concepts That Click~|~YouTube: \href{https://www.youtube.com/@concepts.thatclick}{@concepts.thatclick}}}
  \renewcommand{\headrulewidth}{0pt}
  \renewcommand{\footrulewidth}{0pt}
  \fancypagestyle{plain}{%
    \fancyhf{}%
    \fancyfoot[C]{\footnotesize\textcolor{neutral}{Concepts That Click~|~YouTube: \href{https://www.youtube.com/@concepts.thatclick}{@concepts.thatclick}}}%
    \renewcommand{\headrulewidth}{0pt}%
    \renewcommand{\footrulewidth}{0pt}%
  }
}{}
\@ifclassloaded{beamer}{%
  \setbeamertemplate{background}{%
    \begin{tikzpicture}[remember picture, overlay]
      \node[rotate=45, opacity=0.055, align=center]
        at (current page.center)
        {\Huge\textcolor{neutral}{Concepts That Click}};
    \end{tikzpicture}%
  }
}{}
\makeatother 
\coursecode{JKSSB}
\renewcommand{\thesection}{38.\arabic{section}}
\newtheorem{example}{Example}[section]
\newenvironment{definitionbox}[1]{%
  \par\noindent\textbf{Definition (#1).}\ \itshape
}{\par\vspace{0.3em}}
\title{SMTP, POP and IMAP --- Detailed Lecture Notes}
\author{JKSSB Computer Awareness}
\date{\today}

\begin{document}
\maketitle

\section{Why email needs protocols}
Sending and reading mail involve different programs and different servers.
The sender writes a message in a mail program, that program hands it to a
mail server, the message travels between servers, waits in the recipient's
mailbox, and is finally collected by the recipient's program. For all these
independent programs to work together, they follow fixed rules called
\textbf{protocols}. One protocol sends mail out, and two other protocols
bring it in, each in its own way. Secure transmission and protection against
mail threats sit on top of these three protocols.

\begin{definitionbox}{Email protocol}
A fixed set of rules that a mail program and a mail server follow when
sending, forwarding, or collecting mail, so any standard program can work
with any standard server.
\end{definitionbox}

\begin{center}
\begin{tabular}{p{0.24\textwidth}p{0.66\textwidth}}
\toprule
\textbf{Term} & \textbf{Meaning in this lesson} \\
\midrule
SMTP & Simple Mail Transfer Protocol; the sending protocol. \\
POP3 & Post Office Protocol version 3; a receiving protocol that downloads mail to one machine. \\
IMAP & Internet Message Access Protocol; a receiving protocol that syncs mail across devices. \\
SSL/TLS & Secure Sockets Layer / Transport Layer Security; the secure channel for mail transmission. \\
Spam & Bulk unwanted mail, usually junk or advertising. \\
Phishing & A fake message impersonating a trusted sender to steal secrets. \\
Spoofing & Forging the sender address so mail looks as if it came from someone trusted. \\
\bottomrule
\end{tabular}
\end{center}

The rest of this lesson uses these terms in one consistent sense.
\textbf{SMTP} always means the sending step. \textbf{POP3} always means
downloading to one local machine. \textbf{IMAP} always means keeping every
device in step with the server copy. \textbf{SSL/TLS} always means the
secure transmission channel. \textbf{Spam}, \textbf{phishing}, and
\textbf{spoofing} are three distinct threats defined above.

\section{SMTP: the sending protocol}
\textbf{SMTP} (Simple Mail Transfer Protocol) is the protocol used to send
mail. When the sender presses send, the mail program hands the message to
the sender's mail server using SMTP, and the servers forward it toward the
recipient's mail server using SMTP. SMTP moves the message out of the
sender and along the path between servers until it reaches the recipient's
mailbox.

\begin{definitionbox}{SMTP}
The protocol that sends mail from the sender's program to the mail server
and forwards it between mail servers toward the recipient's mailbox.
\end{definitionbox}

Two limits matter for the exam. First, SMTP only \emph{sends}; it does not
deliver mail into the recipient's hands for reading. Second, the receiving
step that follows is the job of POP3 or IMAP, described next. Any option
that says SMTP downloads mail to the inbox, or that POP3 sends mail,
reverses the direction and is wrong.

\section{POP3: download to one machine}
\textbf{POP3} (Post Office Protocol version 3) is a receiving protocol.
It pulls mail from the mail server down to one local machine. After the
download, the message is typically removed from the server, so the only
copy lives on that one computer. Mail downloaded this way can be read
offline, but a second device will not see the same mailbox, because the
server copy is gone.

\begin{definitionbox}{POP3}
A receiving protocol that downloads mail from the server to one local
machine, usually removing it from the server.
\end{definitionbox}

POP3 suits a single computer where mail is collected and read in one place.
Its weakness is exactly its method: because the mail leaves the server,
bookmarks such as read marks, folders, and deletions do not travel to other
devices.

\begin{example}
A clerk collects office mail only on the office desktop. With POP3, each
morning's mail downloads to that desktop and is removed from the server.
The clerk can read it offline at the desk, but the same messages do not
appear on the clerk's phone, because the server copy is gone.
\end{example}

\section{IMAP: sync across devices}
\textbf{IMAP} (Internet Message Access Protocol) is also a receiving
protocol, but it works differently from POP3. The mail stays on the
server, and the mail program shows a synced view of it: folders, read and
unread marks, deletions, and moves are all mirrored between the device and
the server. When the reader acts on one device, every other device shows
the same change, because each of them reflects the same server copy.

\begin{definitionbox}{IMAP}
A receiving protocol that keeps mail on the server and syncs each device's
view --- folders, read marks, deletions --- with that server copy.
\end{definitionbox}

IMAP suits a reader with several devices, such as a phone and a desktop,
who must see the same inbox everywhere. Its cost is that the mailbox lives
on the server, so reading new mail needs a connection, and the server
mailbox must have room.

\section{POP3 vs IMAP and the full mail flow}
Both POP3 and IMAP receive mail; neither sends it. The difference is the
method. POP3 downloads to one machine and usually clears the server; IMAP
keeps the server copy and syncs every device with it.

\begin{center}
\begin{tabular}{p{0.20\textwidth}p{0.34\textwidth}p{0.36\textwidth}}
\toprule
\textbf{Point} & \textbf{POP3} & \textbf{IMAP} \\
\midrule
Job & Receives mail & Receives mail \\
Method & Downloads to one machine & Syncs a view of the server copy \\
Server copy & Usually removed & Kept on the server \\
Many devices & Poor fit & Designed for it \\
Offline reading & Natural (copy is local) & Needs a connection for new mail \\
\bottomrule
\end{tabular}
\end{center}

The full journey of one message combines the sending protocol with one of
the receiving protocols. The sender's program sends the message by SMTP to
the sender's mail server; that server forwards it by SMTP to the
recipient's mail server, where it waits in the mailbox; the recipient's
program then collects it with POP3 (download) or IMAP (sync). In one line:
sender, then SMTP, then the recipient server, then POP3 or IMAP, then the
reader. Sending is always SMTP; the receiving side is POP3 or IMAP.

\section{To, Cc, BCC, filters, and signatures}
The \textbf{To} field names the main recipients of a message, and
\textbf{Cc} (carbon copy) names others kept in the loop. Everyone who
receives the message can see both lists. \textbf{BCC} (blind carbon copy)
names recipients who are hidden from each other: nobody on the To or Cc
list sees who is on BCC, and BCC recipients do not see each other. The
sender always knows the full list; BCC hides recipients from each other,
not from the sender.

\textbf{Filters}, also called \textbf{rules}, sort incoming mail
automatically. A filter examines the headers and content of each message,
such as the sender, the subject, or words in the body, and then files,
flags, or deletes matching mail --- for example, moving all newsletters to
one folder. A mail \textbf{signature} is different: it is a fixed block of
sender details, such as name, post, and phone number, added at the end of
outgoing messages. It identifies the sender but does not encrypt or
otherwise protect the message.

\section{Digital signature vs encryption, and secure transmission}
A \textbf{digital signature} proves who sent a message and that it was not
changed on the way: it gives authenticity and integrity. \textbf{Encryption}
does something different: it scrambles the message so that only the
intended recipient can read it, giving confidentiality. A signed message
can still be read by anyone --- a signature is not secrecy. An encrypted
message can still be forwarded --- encryption hides the content but does
not by itself prove who sent it. Confusing the two is a tested trap
(TRAP-15): a signature never encrypts the body.

\begin{definitionbox}{Digital signature vs encryption}
A digital signature provides authenticity (who sent it) and integrity (it
was not altered); encryption provides confidentiality (only the intended
recipient can read it).
\end{definitionbox}

When mail travels between programs and servers, the channel itself can be
protected by \textbf{SSL} (Secure Sockets Layer) and its successor
\textbf{TLS} (Transport Layer Security). Together \textbf{SSL/TLS} provide
the secure channel over which mail is transmitted, so that others on the
path cannot read it. For the exam, the answer is direct: the secure email
transmission protocol is SSL/TLS (PYQ-36). The official Junior Assistant
item behind PYQ-16 separately confirms the POP3, BCC, filter, and
signature-versus-encryption facts used in this lesson.

\begin{example}
An item states, ``A digitally signed email can be read only by its
recipient.'' This is wrong. The signature proves the sender and detects
tampering, but anyone can still read the body. Only encryption makes the
body readable solely by the intended recipient. The wrong option has
assigned encryption's job to the signature.
\end{example}

\section{Spam, phishing, spoofing, and malicious content}
\textbf{Spam} is bulk unwanted mail, usually advertising or junk sent to
many addresses at once. Its defining mark is that it is bulk and unwanted;
it clutters the inbox and may also carry fraud or malware. Spam filters move
suspected junk to a separate folder using the sender's reputation and
content rules.

\textbf{Phishing} is a fake message that pretends to be a trusted sender
--- a bank, an office, or a colleague --- in order to steal passwords or
money. It typically creates urgency (``verify now'', ``account blocked'',
``claim within 24 hours'') and asks the reader to click a link or open an
attachment and hand over secrets. The defence is to check the real sender
address and never enter passwords through a link reached from mail.

\textbf{Spoofing} forges the sender address so that a message looks as if
it came from someone trusted. Phishing messages often use spoofing, but
the two terms are distinct: phishing is the fraud attempt, and spoofing is
the forged-sender trick it may use. A \textbf{malicious link} leads to a
fake login page or a malware download; hovering over it shows an address
that does not match the claimed sender. A \textbf{malicious attachment}
carries malware behind an innocent-looking file name. The safe rule is to
verify an unexpected message through another channel before clicking any
link or opening any attachment.

\section{Exam retrieval and common confusions}
Six distinctions prevent most errors on this topic.
\begin{enumerate}
  \item SMTP \textbf{sends}; POP3 and IMAP \textbf{receive}. Reversing the
    direction is a common wrong answer.
  \item POP3 \textbf{downloads} mail to one local machine (usually clearing
    the server); IMAP \textbf{syncs} every device with the server copy.
    Do not swap them (TRAP-17).
  \item A digital signature proves the sender and detects tampering; it
    does \textbf{not} encrypt the body. Encryption alone provides secrecy
    (TRAP-15).
  \item BCC hides recipients \textbf{from each other}; the sender always
    knows the full list, and To/Cc are visible to all (PYQ-16).
  \item Secure transmission of email is \textbf{SSL/TLS} (PYQ-36). A
    signature is not a transmission protocol.
  \item Spam is \textbf{bulk unwanted} mail; phishing is \textbf{fake
    trusted mail that steals secrets}; spoofing is the \textbf{forged
    sender address} either may use.
\end{enumerate}

\begin{example}
An item asks which protocol suits reading the same office mailbox from a
phone and a desktop. Trace the methods: POP3 downloads to one machine and
clears the server, so the second device misses the mail; IMAP keeps the
server copy and syncs every device with it. The answer is IMAP. If the stem
instead asks which protocol sends the message out, the answer is SMTP, and
if it asks which channel makes transmission secure, the answer is SSL/TLS.
\end{example}

\subsection*{Exercises}
\begin{enumerate}[label=\arabic*.]
  \item State the full form of SMTP, POP3, IMAP, SSL, and TLS, and give the
    job of each of the first three.
  \item Explain the difference between POP3 and IMAP when the same mailbox
    is read from two devices.
  \item Trace one message from the sender to the reader, naming the
    protocol used at each step.
  \item Distinguish a digital signature from encryption, and state which
    protocol secures email transmission.
  \item Define spam, phishing, and spoofing, and state the safe rule for
    unexpected links and attachments.
\end{enumerate}

\subsection*{Solutions}
\begingroup\small
\begin{enumerate}[label=\arabic*.]
  \item SMTP is Simple Mail Transfer Protocol and sends mail; POP3 is Post
    Office Protocol version 3 and downloads received mail to one machine;
    IMAP is Internet Message Access Protocol and syncs received mail across
    devices; SSL is Secure Sockets Layer and TLS is Transport Layer
    Security, together the secure transmission channel.
  \item With POP3 the mail downloads to the first machine and usually
    leaves the server, so the second device does not see it; with IMAP the
    mail stays on the server and both devices show the same synced view of
    folders, read marks, and deletions.
  \item The sender's program sends the message by SMTP to the sender's
    server; the servers forward it by SMTP to the recipient's server, where
    it waits; the recipient collects it with POP3 (download) or IMAP
    (sync).
  \item A digital signature provides authenticity and integrity (who sent
    it, unaltered); encryption provides confidentiality (only the recipient
    can read it). Secure transmission uses SSL/TLS.
  \item Spam is bulk unwanted mail; phishing is a fake trusted message
    that steals secrets through urgency plus a link or attachment;
    spoofing forges the sender address. Verify an unexpected message by
    another channel before clicking any link or opening any attachment.
\end{enumerate}
\endgroup
\end{document}
